\chapter{Consumo Especifico}

O consumo específico de energia elétrica é um indicador utilizado para avaliar o desempenho energético da unidade consumidora em função de sua atividade principal. Esse indicador permite comparar a eficiência relativa entre unidades semelhantes, independentemente de seu porte, auxiliando na identificação de oportunidades de melhoria.

A forma de cálculo do consumo específico pode variar conforme o tipo de uso ou função da edificação:

\begin{itemize}
    \item \textbf{Escolas:} kWh por aluno/mês ou kWh por m²/mês;
    \item \textbf{Unidades de saúde:} kWh por leito/mês ou kWh por atendimento/mês;
    \item \textbf{Unidades administrativas:} kWh por m²/mês.
\end{itemize}

Neste estudo, o consumo específico calculado está apresentado na  Tabela~\ref{tab:consumo-especifico}. 

\vspace{1cm}



\begin{table}[h]
\centering
\begin{tabular}{|>{\columncolor[HTML]{1a3a5c}}p{7cm}|p{8cm}|}
\hline
\textcolor{white}{\textbf{<<parametro>>}} & <<dadoParametro>> \\
\hline
\textcolor{white}{\textbf{Consumo médio mensal}} & <<consMedio>> \\
\hline
\textcolor{white}{\textbf{Consumo específico estimado}} & <<dadoConsEspecifico>> \\
\hline
\end{tabular}
\caption{Cálculo do consumo específico de energia elétrica}
\label{tab:consumo-especifico}
\end{table}







%\begin{table}[H]
 %   \centering
  %  \includegraphics[width=\textwidth]{Figuras/Tabela 4.png}
   % \caption{Cálculo do consumo específico de energia elétrica}
%\label{tab:consumo-especifico}
%\end{table}




